\section{Diskussion, Limitationen und Ausblick}
\label{sec:diskussion}

\subsection{Beantwortung der Forschungsfragen}

\paragraph{Kann Q-Learning Kartenzählinformationen nutzen?}
Ja. Die True-Count-Bucket-Analyse zeigt systematische Unterschiede in
der Aktionswahl. Besonders deutlich ist die um fast 4,8 Prozentpunkte
niedrigere Hit-Rate des Counting-Agenten bei einem True Count von
mindestens drei. Dass sich die relative Performance gegenüber der
Baseline über die Buckets verändert, spricht ebenfalls für eine
aktive Nutzung. Ein strenger Kausalnachweis würde jedoch eine
Ablation benötigen, bei der einzelne Zusatzmerkmale entfernt oder
permutiert werden.

\paragraph{Verbessern die erweiterten Features die Performance?}
Im finalen aggregierten Experiment ja, aber nur geringfügig. Der
Counting-Agent verbessert den durchschnittlichen Reward um
$0{,}001794$ und die Win Rate um 0,087 Prozentpunkte. Dieser Vorteil
tritt in allen drei Trainingsseed-Paaren auf, ist wegen der sehr
kleinen Zahl unabhängiger Trainingsläufe aber nicht robust
signifikant. Dem Gewinn steht eine rund 406-mal größere Q-Tabelle
gegenüber. Für tabellarisches Q-Learning ist das Verhältnis von
zusätzlicher Komplexität zu Performancegewinn ungünstig.

\paragraph{Wie viele Episoden werden bis zur Konvergenz benötigt?}
100 Millionen Episoden reichen aus, damit Counting zur Baseline
aufschließt, nicht aber für einen belastbaren Konvergenznachweis. Die
Checkpoint-Leistung schwankt weiterhin und die Q-Tabelle wächst noch.
Auch die konstante Lernrate und die verbleibende Exploration von
$\epsilon=0{,}1$ erfüllen keine hinreichenden Bedingungen, um
theoretische Konvergenz zu behaupten. Längere Läufe könnten weitere
Erkenntnisse liefern, sollten aber mit zusätzlichen Seeds und
Checkpoint-Evaluationen kombiniert werden.

\subsection{Herausforderungen}

\paragraph{Exploration und großer Zustandsraum.}
Die Baseline besucht ihre kleine Zustandsmenge schnell und kann
Q-Werte häufig aktualisieren. Counting verteilt dieselbe Zahl an
Episoden auf mehr als 100.000 Zustände. Viele Kombinationen treten
selten auf, wodurch einzelne Schätzungen verrauscht bleiben. Ein
längerer $\epsilon$-Decay hilft bei der Abdeckung, verzögert aber die
Ausnutzung guter Policies.

\paragraph{Stochastik und Evaluation.}
Millionen Testepisoden reduzieren das Monte-Carlo-Rauschen einer
festen Policy stark. Sie ersetzen jedoch keine unabhängigen
Trainingsläufe. Die Ergebnisse von Seed 123 demonstrieren, dass ein
Modell trotz identischer Hyperparameter anders enden kann. Künftige
Aussagen über Signifikanz sollten deshalb auf mehr Trainingsseeds
statt ausschließlich auf noch mehr Episoden pro Eval-Seed beruhen.

\paragraph{Rechen- und Speicheraufwand.}
Das finale Training dauerte trotz paralleler Ausführung rund 13
Stunden. Counting-Checkpoints sind ungefähr 6,3 MB groß,
Baseline-Checkpoints nur rund 32 KB. Noch längere Läufe erhöhen
Laufzeit und Artefaktmenge, ohne dass tabellarisches Q-Learning
zwischen ähnlichen Zuständen Wissen übertragen kann.

\subsection{Grenzen der Aussagekraft}

Die Simulationsumgebung bildet nur einen Ausschnitt realer
Blackjack-Regeln ab. Naturals erhalten keinen 3:2-Payout; Double
Down, Split, Insurance, Surrender und variable Einsätze fehlen.
Gerade Kartenzählen wird in der Praxis überwiegend für die Anpassung
von Einsätzen verwendet. Der hier gemessene Nutzen unterschätzt oder
verändert daher potenzielle Effekte eines vollständigen Spiels.

Running Count, True Count und verbleibende Decks sind redundant. Dies
erleichtert dem Agenten möglicherweise bestimmte Entscheidungen,
vergrößert aber den Zustandsraum unnötig. Außerdem wird der True
Count ganzzahlig abgeschnitten und geclippt. Die Bucket-Evaluation
ordnet eine Episode nach ihrem Anfangszustand ein, obwohl sich der
Count während der Runde verändern kann. Die Aktionsraten
berücksichtigen zwar den aktuellen Count, der Episodenreward bleibt
jedoch dem Start-Bucket zugeordnet.

Die vorhandene Basic-Strategy-Visualisierung dient nur als
qualitative Hilfe. Da die Projektumgebung eine reduzierte Regelmenge
verwendet und die Visualisierung nicht alle Zustandsdetails des
Counting-Agenten realistisch rekonstruiert, ist sie kein formaler
Optimalitätsnachweis.

Auch die exakte Reproduzierbarkeit eines vollständigen Neutrainings
ist eingeschränkt. Umgebung und Action Space werden explizit
geseedet, die globale NumPy-Zufallsquelle für die $\epsilon$-greedy
Aktionswahl jedoch nicht pro Trainingsprozess gesetzt. Die
gespeicherten Modelle erlauben eine reproduzierbare greedy Evaluation
der berichteten Ergebnisse; ein erneuter Trainingslauf mit denselben
nominellen Seeds muss dagegen nicht bitidentisch enden. Künftig
sollte jeder Worker einen eigenen, explizit gesetzten
Zufallszahlengenerator verwenden.

\subsection{Verbesserungsmöglichkeiten}

Die unmittelbarsten Verbesserungen betreffen die Validität und
Effizienz des bestehenden Q-Learning-Experiments. Separate Agenten
mit ausschließlich True Count, Running Count oder verbleibenden Decks
würden in einer Ablationsstudie zeigen, welche Merkmale tatsächlich
Nutzen stiften. Auf dieser Grundlage könnte eine kompaktere
Zustandsrepräsentation entwickelt werden, die redundante
Kombinationen vermeidet. Ein längerer Lauf mit bis zu 500 Millionen
Episoden ist erst danach sinnvoll und sollte durch feste Checkpoints
sowie mindestens zehn Trainingsläufe abgesichert werden.

Die Umgebung sollte um Natural-Payout, Double Down, Split, Surrender
und variable Einsätze erweitert werden. Erst dann lässt sich
untersuchen, ob ein positiver Count tatsächlich in einen positiven
erwarteten Gewinn oder eine geeignete Einsatzstrategie übersetzt
werden kann. Eine passende Reward-Funktion könnte statt eines
konstanten Rundenergebnisses den Nettoertrag inklusive Einsatzhöhe abbilden.

Zusätzlich sollte das Evaluationsprotokoll stärker auf
Trainingsunsicherheit ausgerichtet werden. Mehr unabhängige
Trainingsläufe sind aussagekräftiger als eine weitere Vergrößerung
der bereits umfangreichen Evaluation einzelner Modelle. Explizit
gesetzte Zufallszahlengeneratoren pro Worker und unveränderliche
Manifeste für alle ausgewerteten Artefakte würden die
Wiederholbarkeit verbessern.

\subsection{Ausblick}

Über diese direkten Verbesserungen hinaus bieten sich mehrere
methodische Erweiterungen an. Transfer Learning könnte die
Baseline-Policy als Initialisierung oder Lehrer für den
Counting-Agenten verwenden. Dadurch müsste die grundlegende
Hit/Stand-Strategie nicht erneut für jeden Count-Zustand gelernt werden.

Für größere Zustands- und Aktionsräume bietet sich anschließend
Funktionsapproximation an. Ein Deep Q-Network (DQN) ist aufgrund des
diskreten Aktionsraums eine naheliegende wertbasierte Alternative und
könnte Wissen zwischen ähnlichen Zuständen teilen. Proximal Policy
Optimization (PPO) würde als Policy-Gradient-Verfahren einen
methodisch anderen Vergleich ermöglichen und insbesondere bei
zusätzlichen Aktionen relevant werden. Beide Verfahren führen jedoch
neue Hyperparameter, höhere Rechenkosten und mögliche
Trainingsinstabilität ein. Ein fairer Vergleich muss deshalb dasselbe
Seed-, Checkpoint- und Evaluationsprotokoll wie das tabellarische
Q-Learning verwenden.

Multi-Agent Reinforcement Learning ist für die aktuelle Aufgabe
weniger vordringlich, weil der Dealer einer festen Regel folgt und
nicht lernt. Es könnte erst bei mehreren interagierenden Spielern
oder adaptiven Gegenspielern relevant werden. Ein Einsatz in realen
Spielsituationen wäre nur nach einer deutlich realistischeren
Umgebung und einer rechtlichen sowie ethischen Bewertung sinnvoll.
